\section{Exergy: where the work is actually lost}
\label{sec:exergy}

The first law says this store is perfect. At cyclic steady state the PCM
returns to its own initial state, so $\eta_{\mathrm{storage}} \equiv 1$ by
construction: every joule that goes down the well comes back up. That is not
a modelling triumph, it is a statement that the energy balance cannot see
what the store costs. The whole of that cost is a degradation of temperature,
and only the second law prices it.

This section sets out the audit added at v0.13, what it found, and --- just
as importantly --- which parts of it are checks that can fail and which are
bookkeeping identities that cannot.

\subsection{The dead state, and why there are two reservoirs}
\label{sec:exergy:deadstate}

The plant exchanges heat with two quite separate reservoirs, and an earlier
draft of this analysis wrongly treated them as one. The ORC condenser rejects
to the sink at $T_{\mathrm{sink}} = \SI{20}{\celsius}$. The heat-pump
evaporator does \emph{not} draw from the sink: its source is geothermal water
at $T_{\mathrm{source}} = \SI{60}{\celsius}$, produced from wells elsewhere in
the cluster.

We take $T_0 = T_{\mathrm{sink}}$. It is the lowest temperature in the system,
so every exergy in the audit is positive, and the heat dumped at the ORC
condenser is genuinely unrecoverable, which is what a dead state should mean.
The consequence is that the geothermal stream is not free: it enters the
balance as a second input,
%
\begin{equation}
  W_{\mathrm{el,in}} + \mathcal{E}_{\mathrm{geo}}
  \;=\; W_{\mathrm{el,out}} + \mathcal{E}_{\mathrm{reinj}}
        + \sum_j I_j .
  \label{eq:exergy:balance}
\end{equation}

\subsection{The geothermal source is finite, and its flow is derived}
\label{sec:exergy:resource}

Until v0.13 the \SI{60}{\celsius} source was effectively unlimited. The model
drew whatever flow the duty required, no output recorded how much that was,
and nothing objected to a design quietly demanding five times the geothermal
field. The flow is now reported, and converted into a well count.

It is \emph{derived}, not prescribed --- the same choice already made for
$N_{\mathrm{wells}}$:
%
\begin{equation}
  \dot m_{\mathrm{geo}}
  = \frac{\dot Q_{\mathrm{src}}}
         {h(T_{\mathrm{source}}) - h(T_{\mathrm{source}} - \Delta T_{3A4A})},
  \qquad
  N_{\mathrm{geo}} = \frac{\dot m_{\mathrm{geo}}}{\dot m_{\mathrm{geo,well}}} .
  \label{eq:exergy:mgeo}
\end{equation}
%
Prescribing a plant-level flow would put the uncertainty in a quantity nobody
has measured. Deriving it moves the uncertainty into the per-well yield
$\dot m_{\mathrm{geo,well}}$, which is field data. Note also that the source
exergy is homogeneous of degree one in $\dot m_{\mathrm{geo}}$, and
\eqref{eq:exergy:mgeo} makes that flow proportional to the duty; both
conventions below therefore reduce to the duty, the glide and the two
temperatures. Deriving the flow rather than prescribing it costs the
second-law analysis nothing.

The default $\dot m_{\mathrm{geo,well}} = \SI{12.5}{\kilo\gram\per\second}$ is
\emph{provisional}. It is the producer in a \SI{4500}{\metre},
\SI{206}{\celsius} repurposed-well case study, which is a hotter and deeper
well than ours; published yields for repurposed oil and gas wells run from
about \SI{1}{\kilo\gram\per\second} up to that figure, while co-produced water
from depleted fields can exceed it. $N_{\mathrm{geo}}$ is directly
proportional to it and should be read as a scaling result, not a prediction.

\subsection{Two conventions, and why the choice is not cosmetic}
\label{sec:exergy:convention}

What should the plant be charged for?

\begin{description}
  \item[\texttt{stripped}] only the exergy the evaporator removes between
    $T_{\mathrm{source}}$ and the reinjection temperature. Reinjection is
    free.
  \item[\texttt{resource}] everything the well lifts, relative to $T_0$, with
    the reinjected stream booked as a named \emph{loss} attributed to no
    component --- it leaves the boundary intact and is not destroyed inside
    it.
\end{description}

For a dedicated producer, drilled and paid for whether or not its exergy is
used, \texttt{resource} is the honest convention, and it is the default. The
choice matters because the two give \emph{opposite} guidance on
$\Delta T_{3A4A}$, as Table~\ref{tab:exergy:dt3a4a} shows. Under
\texttt{stripped} the second-law efficiency falls monotonically and the
optimiser is pushed toward the smallest possible glide --- which is exactly
the configuration that needs the most geothermal wells. Under
\texttt{resource} there is an interior optimum near \SI{20}{\kelvin}.

\begin{table}[htbp]
  \centering
  \caption{The source-side trade. $\Delta T_{3A4A}$ was a quiet approach
  parameter before v0.13; it is the central design variable of the source
  side. Every row is a full cyclic-steady-state run on the plant basis of
  \S\ref{sec:exergy:basis}, so these are directly comparable with
  Table~\ref{tab:exergy:psi}.}
  \label{tab:exergy:dt3a4a}
  \begin{tabular}{rrrrrrrr}
    \toprule
    $\Delta T_{3A4A}$ & $T_{13h}$ & COP & $\eta_{\mathrm{RTE}}$ &
    $\dot m_{\mathrm{geo}}$ & resource & $\psi_{\mathrm{str}}$ & $\psi$ \\
    (K) & (\si{\celsius}) & & & (rel.) & used & & \\
    \midrule
     5 & 50.0 & 2.582 & 0.3310 & 1.000 & \SI{22.7}{\percent} & 0.2833 & 0.1899 \\
    10 & 45.0 & 2.473 & 0.3172 & 0.486 & \SI{42.6}{\percent} & 0.2764 & 0.2356 \\
    15 & 40.0 & 2.373 & 0.3045 & 0.315 & \SI{59.7}{\percent} & 0.2698 & 0.2505 \\
    20 & 35.0 & 2.280 & 0.2927 & 0.229 & \SI{73.9}{\percent} & 0.2634 & \textbf{0.2544} \\
    25 & 30.0 & 2.194 & 0.2818 & 0.178 & \SI{85.2}{\percent} & 0.2571 & 0.2533 \\
    30 & 25.0 & 2.114 & 0.2716 & 0.144 & \SI{93.3}{\percent} & 0.2511 & 0.2498 \\
    \bottomrule
  \end{tabular}
\end{table}

The design point sits at $\Delta T_{3A4A} = \SI{10}{\kelvin}$, and
Table~\ref{tab:exergy:dt3a4a} says that is badly placed for the resource.
Moving to \SI{20}{\kelvin} costs \SI{7.8}{\percent} of COP and
\SI{7.7}{\percent} of $\eta_{\mathrm{RTE}}$, but cuts the geothermal flow to
\SI{47}{\percent} of its present value --- from \num{8.07} producers to
\num{3.80}, and the whole field from \num{23.7} wells to \num{19.4}. At plant level that
is not a close call. What bounds the move is a constraint the model did not
previously represent at all: the minimum reinjection temperature the
formation will tolerate. \texttt{T\_reinject\_min\_C} now exists for it,
deliberately slack, and warns rather than raising, because the project has no
measured value yet.

\subsection{What can fail, and what cannot}
\label{sec:exergy:gates}

Equation~\eqref{eq:exergy:balance} looks like a closure check, and the
temptation is to print its residual and call a small number a pass. That
would be worthless here, and it is worth saying why.

Each machine bucket \emph{is} a genuine test. The destruction in the heat
pump is summed from the entropy generated inside each of its nine components,
from their own state points and flow fractions, and compared against a
boundary balance built only from stream states. Those are independent
computations, and they agree to $\sim\!10^{-15}$ relative. Reintroducing the
defect of \S\ref{sec:exergy:defect} makes the heat-pump gate fire at
$4.5\times10^{-3}$, which is how we know it can.

A third gate was added at v0.15 and is the only one that spans two functions.
$\eta_{\mathrm{RTE}}$ is computed inside the audit from the scaled budget and
the parasitics, and independently in \texttt{performance\_indices} from the
energy the field actually moved; the two agree to the last bit. It exists
because the manuscript, not the code, caught the defect of
\S\ref{sec:exergy:basis}, and this is the check that would have caught it.

The \emph{global} sum cannot fail. Because $T_{3c} = T_{4c}$ --- the water
leaves the condenser and enters the borehole with no state change between ---
and because the discharge stream is likewise shared, the borehole term, taken
as the difference of the two water-stream exergies, cancels the two exchanger
terms exactly. The identity is then $W_{\mathrm{el,in}} +
\mathcal{E}_{\mathrm{stripped}} - W_{\mathrm{el,out}}$ on both sides. An
earlier draft printed this as a closure check reading exactly
\texttt{0.00e+00} and reported it as a pass. A residual of exactly zero over
twenty terms of order $10^{3}$ is not a triumph, it is a tell.

So the borehole number below is correct but \emph{unverified}. Its independent
value requires the resolved integral over $z$ and $t$, which needs an entropy
constitutive curve $s(E')$ alongside the existing $E'(T)$; that is the next
step and is not in v0.13.

\subsection{Two bases that had to be made one}
\label{sec:exergy:basis}

Until v0.15 this audit and $\eta_{\mathrm{RTE}}$ were computed on different
bases and over different boundaries, and the comparison between them --- the
central claim of \S\ref{sec:exergy:psi} --- was not like for like.

$\eta_{\mathrm{RTE}}$, as defined in \S\ref{sec:cyclic}, uses the energy the
field \emph{actually moved} and charges the borehole pumping: the charge pump
adds to the input, the discharge pump subtracts from the output. The exergy
audit did neither. It evaluated every term on the \emph{target} duties of the
energy budget, and charged no parasitics at all. Quoting
``$\num{0.317}$ becomes $\num{0.246}$'' therefore compared two numbers that
did not share a control volume.

The audit is now rebased. At cyclic steady state the storage efficiency is
identically unity and $\lambda$ is forced to zero, so
$\Delta E_{\mathrm{charge}} = \Delta E_{\mathrm{discharge}}$ and
$\Delta E_{\mathrm{out,HP}} = \Delta E_{\mathrm{in,ORC}}$: both half-cycles
scale by the same factor and the rebase is a single multiplication, here
$\num{1.0446}$. The parasitics are then charged with
$\eta_{\mathrm{RTE}}$'s own sign convention, and because they are work
dissipated as friction in the well they appear \emph{both} as an input and as
a destruction --- \SI{703}{\kilo\watt\hour} per cycle, \SI{2.7}{\percent} of
the total, on its own line in Table~\ref{tab:exergy:components}.

With that done the two share a boundary exactly, which is what the third gate
of \S\ref{sec:exergy:gates} now asserts. The second-law efficiency moves from
$\num{0.246}$ to $\num{0.236}$, and the stripped value from $\num{0.289}$ to
$\num{0.276}$.

It is worth recording how this was found. It was not found by the code, or by
any check in it: it was found while writing the manuscript, by asking what
control volume $\psi$ and $\eta_{\mathrm{RTE}}$ each referred to and
discovering that the answer differed. A definition written out in prose has to
name its boundary; a definition written in code does not.

\subsection{The defect the gate caught}
\label{sec:exergy:defect}

The heat-pump gate did not close on first construction. It sat at
$5\times10^{-3}$, and the cause was the water mass flow rate, set as
$\dot Q / (c_p \Delta T)$ with $c_p$ evaluated at the mean temperature. Over
the \SIrange{105}{160}{\celsius} charging rise $c_p$ varies about
\SI{3}{\percent}, so the flow was \SI{0.096}{\percent} inconsistent with the
model's own enthalpy data. That is invisible to every energy balance in the
model, because both sides of those balances use the same wrong flow; it is
visible to the exergy audit, which fabricated \SI{7.17}{\kilo\watt} out of it.

There was a trap in the repair. The formula appears twice, written
independently. Fixing the copy in \texttt{energy\_budget} alone changed no
reported number at all --- an encouraging no-op that would have been taken as
confirmation, and was nothing of the kind, because the live copy is in
\texttt{run\_cycle} and was untouched. Both are now enthalpy differences. The
effect on reported results is small: $\eta_{\mathrm{RTE}}$ moves
\SI{+0.0044}{\percent}, the largest single move is \SI{-0.27}{\percent} in
pumping power.

A second and larger change was measured and then \emph{declined}. Using the
realised $T_{2d}$ rather than the specified glide to set the discharge flow
cuts the reported deviation from \SI{+4.46}{\percent} to
\SI{+1.45}{\percent} and lifts $\eta_{\mathrm{RTE}}$ by \SI{0.28}{\percent}.
It is not adopted, for two reasons: it changes the model's specification
(``the flow is pinned by the glide'') rather than fixing an approximation,
and it is partly circular, since feeding one pass's outlet into the next
pass's flow definition drives the delivered energy toward the target by
construction without iterating to a fixed point. It is recorded in DN-25 as a
candidate, with its measured effect, rather than smuggled in under cover of
an exergy release.

\subsection{The destruction table}
\label{sec:exergy:table}

Table~\ref{tab:exergy:components} is the audit at the design point, per cycle.

\begin{table}[htbp]
  \centering
  \caption{Exergy destruction by component, one cycle at CSS,
  $T_0 = \SI{20}{\celsius}$, \texttt{resource} convention. Inputs
  \SI{31813}{\kilo\watt\hour} electrical and \SI{11019}{\kilo\watt\hour}
  geothermal; outputs \SI{10000}{\kilo\watt\hour} electrical and
  \SI{6326}{\kilo\watt\hour} reinjected.}
  \label{tab:exergy:components}
  \begin{tabular}{lrr}
    \toprule
    Component & $I$ (kWh) & share \\
    \midrule
    HTHP condenser        & 4584.4 & \SI{17.4}{\percent} \\
    ORC evaporator        & 3825.5 & \SI{14.5}{\percent} \\
    HTHP throttle 7h--8h  & 3034.4 & \SI{11.5}{\percent} \\
    HTHP compressor HP    & 1915.0 & \SI{7.3}{\percent} \\
    borehole\footnotemark[1] & 1740.3 & \SI{6.6}{\percent} \\
    HTHP motor            & 1573.2 & \SI{6.0}{\percent} \\
    ORC condenser         & 1305.7 & \SI{4.9}{\percent} \\
    HTHP evaporator       & 1233.4 & \SI{4.7}{\percent} \\
    HTHP compressor LP    & 1222.1 & \SI{4.6}{\percent} \\
    ORC turbine 2         & 1214.6 & \SI{4.6}{\percent} \\
    HTHP throttle 9h--4h  & 1168.3 & \SI{4.4}{\percent} \\
    ORC regenerator       &  830.2 & \SI{3.1}{\percent} \\
    \textbf{borehole pumping} &  703.0 & \SI{2.7}{\percent} \\
    HTHP IHX-2            &  671.2 & \SI{2.5}{\percent} \\
    ORC generator         &  549.8 & \SI{2.1}{\percent} \\
    ORC turbine 1         &  511.0 & \SI{1.9}{\percent} \\
    HTHP IHX-1            &  320.3 & \SI{1.2}{\percent} \\
    ORC pumps, mixer, separator & 12.2 & \SI{0.0}{\percent} \\
    \midrule
    total                 & 26414.6 & \\
    \bottomrule
  \end{tabular}

  \footnotesize\footnotemark[1]\,By difference, not independently verified;
  see \S\ref{sec:exergy:gates}.
\end{table}

Three readings, two of them not what we expected.

\paragraph{The exchangers beat the machines.} The condenser, the two
evaporators, the ORC condenser and the borehole together account for
\SI{48.0}{\percent} of all destruction. The compressors and turbines account for
\SI{18.5}{\percent}. This is a heat-transfer plant, not a turbomachinery plant
--- which is convenient, because heat transfer is what this model resolves and
turbomachinery is what it stipulates.

\paragraph{A valve destroys more than the high-pressure compressor.} The
throttle into the separator, at \SI{11.5}{\percent}, is the third-largest
single item and exceeds either compressor. That is a design signal rather than
a modelling artefact: an expander, or a different flash arrangement, is worth
more here than any plausible improvement in compressor efficiency.

\paragraph{The borehole is only \SI{6.6}{\percent}.} It is fourth on the list,
and it is the part of the plant nobody else has modelled, so it earns its
place. But the honest reading is that the store is not where this system
bleeds. The case for the cascade and for tight approach temperatures has to be
made on well count and on hardware, not on lost work.

\subsection{What the second law says about the headline number}
\label{sec:exergy:psi}

\begin{table}[htbp]
  \centering
  \caption{Three efficiencies at the design point. The spread is the result.}
  \label{tab:exergy:psi}
  \begin{tabular}{llr}
    \toprule
    $\eta_{\mathrm{RTE}} = W_{\mathrm{out}}/W_{\mathrm{in}}$
      & source treated as free & 0.3172 \\
    $\psi_{\mathrm{stripped}}$
      & source priced at what was taken & 0.2764 \\
    $\psi$
      & source priced at what the well delivers & 0.2356 \\
    \bottomrule
  \end{tabular}
\end{table}

$\eta_{\mathrm{RTE}}$ books the \SI{60}{\celsius} heat as costless. It is not.
Charging only the exergy the evaporator strips already takes the efficiency
down to \num{0.276}; charging the whole resource the well lifts takes it to
\num{0.236}. A referee will ask about this, so it is better that we publish
it.

The most uncomfortable number in the audit is not in
Table~\ref{tab:exergy:components} at all. The reinjected stream carries
\SI{6326}{\kilo\watt\hour} of exergy out of the plant every cycle, against an
net electrical output of \SI{10092}{\kilo\watt\hour}. We discard
\SI{57.4}{\percent} of the geothermal exergy at the wellhead --- more than we
use. Table~\ref{tab:exergy:dt3a4a} says the remedy is not hotter wells but
deeper cooling of the wells we have, and that the obstacle is the reinjection
limit, which we cannot yet quantify.

\subsection{Wells, which is the number an operator asks for}
\label{sec:exergy:wells}

Because the geothermal side is now counted, the model reports the whole field
for the first time: \num{15.62} storage wells and \num{8.07} geothermal
producers, \num{23.69} in total, or \num{22.7} wells per MW\textsubscript{e}.
Until v0.13 we reported \num{15.62} and called it the well count. The
geothermal side is not a minor appendage: it is a third of the boreholes.

One caution carries over from \S\ref{sec:exergy:resource}. Because
$\dot m_{\mathrm{geo}}$ is derived rather than capped, nothing now stops a
parametric study from demanding an arbitrarily large geothermal flow to buy
COP. $N_{\mathrm{geo}}$ must therefore be carried as a reported indicator with
the same standing as $N_{\mathrm{wells}}$, not as a diagnostic. A factorial
study that ranks designs on $\eta_{\mathrm{RTE}}$ alone will find the
free-source corner and recommend it.

\subsection{Two indicators added with the audit}
\label{sec:exergy:kpi}

Two things the exergy work made it obvious we were not reporting.

\paragraph{The four pinches, individually.} \texttt{exchanger\_UA} has
computed all four since v0.11, but only their minimum was exported. Which
exchanger is tightest is design information, and a design can move the
binding one without the minimum changing at all:

\begin{center}
\begin{tabular}{@{}lrrrr@{}}
\toprule
 & HTHP evap & HTHP cond & ORC evap & ORC cond \\
\midrule
pinch [\si{\kelvin}] & 5.00 & \textbf{10.00} & 5.01 & 5.00 \\
$UA$ [\si{\kilo\watt\per\kelvin}] & 463.1 & 241.9 & 413.2 & 1135.1 \\
\bottomrule
\end{tabular}
\end{center}

Three of the four sit hard against the pinch they were given, so
$\Delta T_{\mathrm{pinch,ORC}}$, $\Delta T_{\mathrm{pinch,HPE}}$ and
$\Delta T_{E,\mathrm{sink}}$ are all binding. $\Delta T_{2H3C}$ is the one
loose approach in the plant, and the only one with room to be tightened
without an immediate area penalty.

\paragraph{Power density beside energy density.} $\rho_E$ says how much
energy a cubic metre of store holds and nothing about how fast it can be
delivered; the two are traded against each other by the discharge window. At
the design point

\[
  \rho_E = \SI{170.9}{\kilo\watt\hour\per\meter\cubed},
  \qquad
  \rho_P = \SI{2.415}{\kilo\watt\per\meter\cubed}
  \;=\; \SI{0.0439}{\kilo\watt\per\meter}\ \text{of borehole,}
\]

the last figure being \SI{44}{\watt} per metre of well. Reported per metre as
well as per cubic metre because drilling is priced by the metre, and because
it is the number that invites comparison with a conventional borehole heat
exchanger. A store with twice the energy density and half the rate is a
different machine, not a better one, and $\rho_E$ alone cannot tell the two
apart.

\subsection{Two indicators built, measured and removed}
\label{sec:exergy:removed}

Both were implemented and run before being taken out of the reported set. The
findings are kept here because they are worth having on record.

\paragraph{The residual, decomposed.} The cycled fraction $0.9595$ splits into
a \emph{floor} of $0.0109$ --- PCM still molten when discharge ends, which
never returns its latent heat --- and a \emph{headroom} of $0.0296$ --- PCM
never melted at the end of charge, which never stores at all. The headroom is
nearly three times the floor, so the unused material is at the cold end of
the cascade, and the remedy is charge duration or inlet temperature rather
than anything on the discharge side. Removed to keep the indicator set small
for the factorial study; the cycled fraction and the two means it is built
from remain, so the split can be recovered by hand.

\paragraph{A Lorenz reference.} Every stream in this plant glides, so a
Carnot reference would compare the real machine against an ideal one handed
an easier problem; the thermodynamic mean temperature is the correct bound.
On that basis the heat pump achieves $2.4731$ against $5.2667$ reversible, a
ratio of $0.470$, and the ORC $0.1487$ against $0.2527$, a ratio of $0.589$.
\emph{The heat pump sits further from its own ideal than the ORC}, despite
both carrying \num{0.85} isentropic machines --- consistent with its two
throttles, which are \SI{16.3}{\percent} of all destruction and have no ORC
counterpart.

There is no useful \emph{plant-level} Carnot bound to report alongside these:
treated as a pure electricity store, the reversible round trip is exactly
unity, so $\eta_{\mathrm{RTE}}$ divided by it is $\eta_{\mathrm{RTE}}$ again.

It was removed for two reasons. The ratios carry no information
Table~\ref{tab:exergy:components} does not already have --- what they add is
a normalisation, not evidence. And the guard written for them could not be
made to work: at ratios of $0.47$ and $0.59$ the design is far enough from
reversible that a wrong reference still satisfies $\psi\le1$, and four
deliberate corruptions of the reference all failed to trip it. The reason is
structural and worth stating, because it applies well beyond this function:
\textbf{a derived diagnostic cannot be guarded by a bound computed from its
own inputs}. The gates of \S\ref{sec:exergy:gates} work precisely because
their two sides are built from different things.
